% Generated draft: requires mathematical and source review.
\documentclass[11pt]{article}
\usepackage[margin=25mm]{geometry}
\usepackage{fontspec}
\setmainfont{Latin Modern Roman}
\usepackage{amsmath,amssymb,mathtools,unicode-math}
\setmathfont{Latin Modern Math}
\usepackage{enumitem,needspace,xurl,hyperref}
\setlength{\parindent}{0pt}
\setlength{\parskip}{5pt}
\setlength{\emergencystretch}{3em}
\newcommand{\N}{\mathbb{N}}
\newcommand{\Z}{\mathbb{Z}}
\newcommand{\Q}{\mathbb{Q}}
\newcommand{\R}{\mathbb{R}}
\newcommand{\C}{\mathbb{C}}
\newcommand{\F}{\mathbb{F}}
\newcommand{\A}{\mathbb{A}}
\newcommand{\PP}{\mathbb{P}}
\newcommand{\E}{\mathbb{E}}
\newcommand{\eps}{\varepsilon}
\newcommand{\dd}{\,\mathrm{d}}
\begin{document}
\section*{UnsolvedMath ID: 800819}
\textbf{Source problem code:} catalogue entry
\subsection{Definitions and mathematical statement}
Let \(M\) denote the spatial domain appearing in the source, let \(I\subseteq\R\) be a time interval, and let \(u(t,\cdot)\) denote the velocity field of a solution of the Euler equations on \(M\), defined for \(t\in I\). The supplied record does not specify \(M\), the version of the Euler equations, boundary conditions, or the solution class; these remain unresolved parameters of the question.

\textbf{Aging Problem.} For an arbitrary such solution \(u\) and any two distinct times \(t_1,t_2\in I\), can one determine, from the two velocity fields \(u(t_1,\cdot)\) and \(u(t_2,\cdot)\) alone, which field corresponds to the earlier time?

One possible precise interpretation of this question would ask for a single rule
\[
A:\mathcal V\times\mathcal V\longrightarrow\{-1,1\},
\qquad
A\bigl(u(t_1,\cdot),u(t_2,\cdot)\bigr)
=\operatorname{sgn}(t_2-t_1)
\]
for every admissible solution \(u\) and all distinct \(t_1,t_2\) in its interval of existence, where \(\mathcal V\) is the class of admissible instantaneous velocity fields. This is only a candidate interpretation: the source does not define \(\mathcal V\), require a universal rule explicitly, or specify how coincident snapshots and exceptional solutions are to be treated.

The source also asks whether fluid ``aging'' can be quantified. Its stated motivation is that numerical experiments suggest that, starting from a smooth initial velocity \(u_0=u(0,\cdot)\), the velocity becomes ``wrinkled'' and its derivatives generally grow. The record attributes the expression ``regularity deterioration'' for this phenomenon to V. Yudovich. These observations are supplied motivation, not a stated theorem or a specified monotonicity property. No derivative order, norm, aging functional, time direction, growth rate, or quantifier making ``generally'' precise is supplied. Consequently, the request to quantify aging remains an open-ended formulation requiring further hypotheses.
\subsection{Short English statement}
The source asks whether two velocity snapshots from an arbitrary Euler flow reveal their chronological order, and whether the suggested growth of spatial derivatives can quantify fluid ``aging.'' The domain, solution class, exceptions, and intended quantitative criterion are unspecified. The supplied registry labels the problem open; this draft does not independently verify that status.
\subsection{Sources}
\begin{itemize}
\item Local UnsolvedMath registry, permanent ID \(800819\), source problem code \texttt{catalogue entry}, title: \emph{Geometric Hydrodynamics in Open Problems --- Problem 18 (Aging Problem)}. The supplied record attributes the problem to Khesin, Misiolek, and Modin, with proposed year 2022, and labels its status \emph{open}.
\item Khesin, Misiolek, and Modin, \emph{Geometric Hydrodynamics in Open Problems} (2022), Problem 18. Supplied source link: \url{https://arxiv.org/abs/2205.01143}. The background records an access date of 2026-10-16 and extraction method \texttt{arxiv\_latex}; these are supplied provenance details, not a claim of independent access or verification.
\item The supplied statement attributes the term ``regularity deterioration'' to V. Yudovich and reports numerical experiments suggesting aging. No separate bibliographic reference for either attribution or the experiments is supplied.
\end{itemize}
\end{document}
